\section{Cohomology of the schematic Rees family}

\begin{lemma}[Homology commutes with direct sums of modules]
\label{lem:fc-formal-functions-direct-sum}
\lean{EnvelopeProperty.FormalFunctionsDirectSum.homologyIso}
\leanok
Let $R$ be a ring and $(A_d\xrightarrow{f_d}B_d\xrightarrow{g_d}C_d)_{d\geq0}$
a family of complexes of $R$-modules. The natural map
\[
 \bigoplus_d \ker(g_d)/\operatorname{im}(f_d)
 \longrightarrow
 \ker(\bigoplus_d g_d)/\operatorname{im}(\bigoplus_d f_d)
\]
is an isomorphism.
\end{lemma}
\begin{proof}
A direct-sum element is killed by $\bigoplus_d g_d$ exactly when each of its
finitely many nonzero coordinates is killed by $g_d$. Thus its kernel is the
direct sum of the individual kernels. Its image is the direct sum of the
individual images: one inclusion follows by evaluating every coordinate;
for the reverse inclusion choose a preimage in each of the finitely many
nonzero coordinates and take their direct sum. Coordinatewise quotient maps
are surjective, and their direct-sum kernel is exactly the direct sum of the
individual images. Passing to quotients gives the stated isomorphism. These
maps are the maps induced by the original inclusions of summands.
\end{proof}

\begin{lemma}[Cohomology decomposition of the actual Rees Cech complex]
\label{lem:fc-rees-cech-degree-homology}
\uses{lem:fc-formal-functions-direct-sum, lem:fc-rees-cech-degree-complex}
\lean{EnvelopeProperty.FcReesCechDegreeHomology.orderedHomologyIso}
\leanok
Let $X$ be separated over $\operatorname{Spec}R$, let $I\subset R$, and fix
a finite ordered family of affine opens of $X$. For the actual schematic
Rees family, the ordered Cech complex of its structure sheaf, restricted to
$R$, has cohomology canonically isomorphic to the direct sum of the
cohomologies of the original ideal-power degree complexes.
\end{lemma}
\begin{proof}
On each original affine intersection the actual Rees-family sections are
$\bigoplus_{d\geq0}(I\mathcal O_X)^d$. The actual restriction maps preserve
each degree and restrict to the original ideal-power restriction maps.
Because the original affine family is finite, taking the Cech term commutes
with this direct sum. The resulting cochain isomorphism commutes with every
alternating differential. Apply the preceding lemma to the incoming and
outgoing differentials at each cohomological degree.
\end{proof}

\begin{remark}
The direct-sum comparison is $R$-linear. Finite generation over the Rees
algebra also requires retaining its genuine action, which sends degree $d$
to degree $d+e$ under multiplication by an element of $I^e$. Proper
cohomology finiteness and this action are required before applying stable
filtrations or algebraizing compatible infinitesimal sections. Neither is
an assumption of the eventual formal-functions theorem.
\end{remark}

\begin{lemma}[Actual homogeneous Rees multiplication]
\label{lem:fc-formal-functions-rees-multiplication}
\lean{EnvelopeProperty.FormalFunctionsReesMultiplication.powerEquiv_monomial_mul, EnvelopeProperty.FormalFunctionsReesMultiplication.powerRestriction_multiply}
\leanok
Let $R\to A$ be a homomorphism of commutative rings, $I\subset R$, and
$J=IA$. For $a\in I^e$, multiplication by its image defines the original
$R$-linear map $J^d\to J^{e+d}$. Under the genuine Rees coefficient
isomorphism, multiplication of the actual monomials $aT^e$ and $xT^d$
is exactly inclusion of $ax$ in degree $e+d$. For an $R$-algebra map
$A\to B$, these maps commute with the actual ideal-power restriction.
\end{lemma}
\begin{proof}
The image of $a$ belongs to $J^e$ because ideal extension commutes with
powers. Multiplication of this element with $x\in J^d$ belongs to
$J^eJ^d=J^{e+d}$, and distributivity and commutativity prove $R$-linearity.
Multiplication of the actual polynomials gives
$(aT^e)(xT^d)=(ax)T^{e+d}$; its only nonzero coefficient is $ax$ in the
stated degree. Applying the algebra homomorphism to $ax$ gives the
product of its two original images, and its scalar compatibility retains
the same original image of $a$. This proves the full restriction square.
\end{proof}

\begin{lemma}[Original homogeneous multiplication on ideal-power Cech cohomology]
\label{lem:fc-rees-cech-homogeneous-multiplication}
\uses{lem:fc-formal-functions-rees-multiplication, lem:fc-rees-cech-degree-complex}
\lean{EnvelopeProperty.FcReesCechHomogeneousMultiplication.degreeMultiply, EnvelopeProperty.FcReesCechHomogeneousMultiplication.homologyMultiply}
\leanok
For $a\in I^e$, literal multiplication by the image of $a$ on every
original affine intersection defines a cochain map from the genuine
degree-$d$ ideal-power complex to the degree-$(e+d)$ complex, and hence
an $R$-linear map on each of their cohomologies.
\end{lemma}
\begin{proof}
At every ordered intersection multiply a section by the restriction of $a$.
The preceding restriction square shows that each face restriction commutes
with this multiplication. Therefore each finite alternating sum of face
restrictions commutes with multiplication. This gives the cochain map.
It carries cycles to cycles and boundaries to boundaries, so its induced
map on the original quotient of cycles by boundaries is well defined and
$R$-linear. Its action at every intersection is the actual multiplication
of homogeneous Rees polynomials established in the preceding lemma.
\end{proof}

\begin{lemma}[Summand inclusions in the homology comparison]
\label{lem:fc-formal-functions-direct-sum-inclusions}
\uses{lem:fc-formal-functions-direct-sum}
\lean{EnvelopeProperty.FormalFunctionsDirectSumNaturality.inclusion_homology}
\leanok
For every family of module short complexes, the canonical homology
isomorphism from the direct-sum complex to the direct sum of homologies
carries the homology map of each original summand inclusion to the
corresponding direct-sum inclusion of its original homology.
\end{lemma}
\begin{proof}
Include a cycle in one original summand. Its image is the same cycle
in that coordinate and zero in all other coordinates. The coordinatewise
quotient map therefore gives its original homology class in that summand
and zero elsewhere. Every homology class is represented by a cycle, so
this proves the claimed equality on every original homology class.
\end{proof}

\begin{lemma}[Projectivity of the actual schematic Rees family]
\label{lem:fc-formal-functions-projective-family}
\uses{lem:fc-rees-formal-family, lem:fc-rees-formal-restriction}
\lean{EnvelopeProperty.FormalFunctionsProjectiveFamily.family_isProjectiveFactorization}
\leanok
If $X\to\operatorname{Spec}R$ has an actual closed embedding into a finite
dimensional polynomial projective space over $R$, its actual schematic
Rees family has a closed embedding into the same dimensional polynomial
projective space over $\operatorname{Rees}_R(I)$, over that same coefficient
spectrum.
\end{lemma}
\begin{proof}
Base change the original closed embedding along
$\operatorname{Spec}\operatorname{Rees}_R(I)\to\operatorname{Spec}R$.
The resulting map is a closed embedding because it is an actual pullback
of the original one. The genuine polynomial-projective base-change
isomorphism identifies its target with polynomial projective space over
$\operatorname{Rees}_R(I)$, preserving the entire coefficient projection.
Compose the resulting closed embedding with the actual closed schematic
image inclusion defining the Rees family. Both the closed-embedding
property and the coefficient projection identity are preserved by
composition. This constructs the required projective factorization.
\end{proof}

\begin{lemma}[Finite cohomology of the actual Rees family]
\label{lem:fc-formal-functions-rees-cohomology-finite}
\uses{lem:fc-formal-functions-projective-family, lem:fc-rees-cech-chart-sections}
\lean{EnvelopeProperty.FormalFunctionsReesCohomologyFinite.orderedReesCech_finite, EnvelopeProperty.FormalFunctionsReesCohomologyFinite.chartReesCech_finite}
\leanok
Suppose $R$ is Noetherian and $X\to\operatorname{Spec}R$ has an actual
finite-dimensional closed projective embedding. For every finite ordered
affine open cover of $X$, every cohomology module of the actual ordered
structure-sheaf Cech complex of its schematic Rees family is finite over
$\Gamma(\operatorname{Spec}\operatorname{Rees}_R(I),\mathcal O)$.
\end{lemma}
\begin{proof}
The Rees algebra is Noetherian, and the preceding lemma constructs the
actual family projective embedding. Its structure module is finite type
and quasicoherent. The projection from the actual family to $X$ is affine,
so the inverse images of the given affine opens are affine and cover the
whole actual family. Projective coherent-cohomology finiteness on the
standard coordinate cover, followed by the actual affine-cover Cech
comparison, gives finite native Cech cohomology in each positive degree
on this same pulled-back cover. The native and ordered complexes are
homotopy equivalent, so their original cohomology modules are isomorphic.
In degree zero the sheaf gluing comparison identifies ordered Cech
cohomology with global sections, finite by projective coherent-cohomology
finiteness. Finally the actual Rees chart complex isomorphism preserves
the original coefficient action and transports these finite modules.
\end{proof}

\begin{lemma}[Original ideal-power summands on actual Rees cohomology]
\label{lem:fc-rees-cech-degree-inclusions}
\uses{lem:fc-rees-cech-degree-homology, lem:fc-formal-functions-direct-sum-inclusions}
\lean{EnvelopeProperty.FcReesCechDegreeInclusions.homologyInclusion}
\leanok
The literal inclusion of degree-$d$ ideal-power cochains in the actual
direct-sum Rees degree complex is a cochain map. On cohomology, the genuine
Rees degree decomposition carries this map to the literal inclusion of
the original degree-$d$ ideal-power cohomology summand.
\end{lemma}
\begin{proof}
The whole Rees degree differential acts coordinatewise by the genuine
ideal-power differential. It therefore commutes with inclusion of any
one degree, including its zero coordinates in every other degree. This
constructs the cochain map. Its associated homology three-term map is
the literal summand inclusion in the direct sum of the original homology
three-term complexes. The preceding direct-sum homology inclusion theorem
therefore identifies its induced homology map with the stated original
summand inclusion.
\end{proof}

\begin{lemma}[The actual finite Rees cohomology module]
\label{lem:fc-rees-cech-homology-module}
\uses{lem:fc-rees-cech-degree-homology, lem:fc-formal-functions-rees-cohomology-finite}
\lean{EnvelopeProperty.FcReesCechHomologyModule.finiteReesHomology, EnvelopeProperty.FcReesCechHomologyModule.reesModule_smul, EnvelopeProperty.FcReesCechHomologyModule.reesModule_originalScalar, EnvelopeProperty.FcReesCechHomologyModule.originalScalarTower}
\leanok
The direct sum of original ideal-power cohomology modules carries the
actual Rees coefficient action transported through the genuine original
cohomology comparison. For a Noetherian original base, an original closed
projective embedding and an original finite affine cover, this is a finite
module over the actual Rees algebra. Its scalar action on every class is
precisely the original Rees scalar action conjugated by that comparison.
\end{lemma}
\begin{proof}
Restriction of scalars preserves kernels and cokernels of module maps,
so the genuine original Rees Cech homology restricted to $R$ identifies
with the cohomology of the restricted original complex. Compose this
isomorphism with the genuine degree homology decomposition. Transfer the
original Rees coefficient action along the resulting additive equivalence;
the additive groups and original $R$-module structures are retained.
The equivalence is linear for the transferred action by its definition.
The preceding actual finite-cohomology theorem therefore makes this
module finite over the original Rees-base global functions. The actual
affine section ring isomorphism identifies those global functions with
the actual Rees ring, so restriction along its inverse preserves finite
generation. This gives the stated finite Rees module and its exact action
on every original class. Compatibility of this action with polynomial
degrees is a further statement, not an assumption in this construction.
\end{proof}

\begin{lemma}[Actual coherent ideal-power Cech comparison]
\label{lem:fc-rees-cech-ideal-power-complex}
\uses{lem:fc-rees-cech-ideal-power-sections, lem:fc-rees-cech-degree-complex}
\notready
The actual ordered Cech complex of the original coherent ideal-power
module sheaf, restricted to the original coefficient ring, is isomorphic
to the genuine original ideal-power degree complex. Each original affine
intersection, original face restriction and alternating differential is
preserved.
\end{lemma}
\begin{proof}
Every ordered affine intersection is affine because the original scheme
is separated. Apply the preceding whole section equivalence at each
intersection and take products. Naturality identifies every ordered
face restriction with the same original chart ideal-power restriction.
Taking finite alternating sums identifies the differentials. The termwise
isomorphisms therefore form a cochain isomorphism with the stated
original terms and differential.
\end{proof}
